% =====================================================================
\chapter{Cumplimiento de los objetivos de las guías}\label{cap:cumplimiento}

\begin{center}\small
\begin{longtable}{>{\raggedright}p{6.4cm} >{\raggedright}p{1.8cm} >{\raggedright\arraybackslash}p{6.8cm}}
\toprule
\textbf{Objetivo} & \textbf{Estado} & \textbf{Evidencia o lo que falta}\\
\midrule
\endhead
\multicolumn{3}{@{}l}{\textbf{Laboratorio 1: auditoría}}\\
Clasificar variables PV, MV, DV & \cumple & capítulo \ref{cap:planta}, corregido con la planta en marcha\\
Ficha técnica del hardware & \cumple & capítulos \ref{cap:planta} y \ref{cap:drivers}\\
Ruptura de causalidad, lazo abierto & \cumple & sección del capítulo \ref{cap:planta}\\
Arquitectura de 3 capas & \cumple & capítulo \ref{cap:aplicaciones}\\
\multicolumn{3}{@{}l}{\textbf{Laboratorio 2: planos e interfaces}}\\
Plano eléctrico versión 2.0 & \cumple & \codigo{documentacion/planos\_isa/electrico/brazo\_electrico\_v2.dwg} (IEC 60617, 2 hojas, referencias cruzadas); quedan puntos SUPUESTO por verificar en el tablero\\
P\&ID con tags ISA-5.1 & \cumple & \codigo{documentacion/planos\_isa/pid/brazo\_pid\_v2.dwg}, lazos 100 a 104\\
Mockups HMI y dashboard & \parcial & mockups del equipo en \codigo{database/mockups/}; la app y la página son la interfaz real\\
\multicolumn{3}{@{}l}{\textbf{Laboratorio 3: Git y modelo de datos}}\\
Repositorio estructurado, ramas y pull requests & \cumple & estructura según la guía 3, ramas \codigo{feature/*} y Pull Request \#1 en revisión\\
Modelo E-R normalizado en 3FN & \cumple & capítulo \ref{cap:datos}\\
Scripts en \codigo{/database/scripts/} & \cumple & \codigo{brazo\_4gdl.sql} ejecutado en PostgreSQL 16.2, con registro de ejecución\\
\multicolumn{3}{@{}l}{\textbf{Fase 2: modelado y backend}}\\
Modelado y linealización del brazo & \cumple & capítulo \ref{cap:cinematica}: jacobiano y linealización con errores medidos\\
Redes de Petri & \cumple & sección \ref{sec:petri}\\
Triggers contra datos fuera de rango & \parcial & \codigo{verificar\_limites\_variables()} descarta el dato, registra la alerta y saca la planta de operación; probado en PostgreSQL 16.2\\
Gemelo validado en simulación & \cumple & capítulo \ref{cap:mediciones}\\
\multicolumn{3}{@{}l}{\textbf{Fase 3: implementación}}\\
Programa del PLC con PTO y comunicación & \cumple & capítulo \ref{cap:gemeloplc}\\
Comunicación PLC--PC por tramas binarias & \cumple & capítulo \ref{cap:ppi} (PPI en lugar de Modbus, justificado)\\
Middleware en Python & \cumple & capítulo \ref{cap:aplicaciones}\\
Control de acceso por roles & \cumple & sección \ref{sec:acceso}\\
Lazo cerrado en el hardware & \falta & sin encoders: el brazo sigue en lazo abierto con supervisión discreta\\
Telemetría en PostgreSQL y dashboard & \falta & diseño listo; falta conectar la app a la base\\
\bottomrule
\end{longtable}
\end{center}

\section{Siguientes pasos}
\begin{enumerate}
\item Instalar PostgreSQL en la PC del laboratorio y agregar a la app el registro de telemetría.
\item Implementar rampas en el PTO y subir la corriente de X e Y.
\item Medir el error de posición de la garra y verificar los SUPUESTOS.
\item Agregar encoders (o el contador rápido HSC del PLC) para cerrar el lazo de posición.
\end{enumerate}
